% ============ 05. S1-S3 ============
\section{Sustentación oral y requisitos de entrega (S1--S3)}

\subsection{S1 · Dominio del problema y defensa (28\,\%, individual)}

\textbf{Nivel estimado hoy: En desarrollo (3,0)} si se sustenta sin preparación.
El sustrato es Competente--Ejemplar (los materiales permiten responder sobre
cualquier sección), pero nada garantiza que \emph{cada integrante} pueda
sostener el conjunto. Con el techo \emph{nota $\leq$ S1 $+$ 1,0}, un S1 en 3,0
capa la nota individual en 4,0 aunque el documento sea 5,0.

\textbf{Acción:} preparar un banco de preguntas de defensa con respuestas
propias (ya se propuso un banco de 10 en el informe del agente S, con
respuestas-plantilla extraídas de los materiales) y practicar en parejas; cada
integrante debe poder explicar \emph{por qué} se eligió el problema, qué
alternativa se descartó y qué literatura sostiene cada decisión.

\subsection{S2 · Claridad y estructura de la exposición (18\,\%)}

\textbf{Nivel estimado: sustrato alto; requiere construir el guion.} La rúbrica
pide: en 10 minutos, problema $\rightarrow$ qué se sabe $\rightarrow$ objetivos
$\rightarrow$ por qué cabe en el semestre, con diapositivas que apoyen (no que
se lean). Las figuras ya existen (artifacts del proyecto); falta el guion de 10
minutos y ~9 diapositivas del \emph{anteproyecto} (no de la auditoría).

\subsection{S3 · Trabajo en equipo, ética y uso declarado de IA (14\,\%)}

\textbf{Nivel hoy: Insuficiente / No pertinente} — \textbf{vacío crítico}. No
existe la \emph{``Declaración de uso de IA y reparto del trabajo''} (media
página) que exige la rúbrica, y los materiales evidencian uso intensivo de
agentes de IA (auditoría de código, validaciones independientes, redacción,
búsqueda de literatura): no declararla, o declararla en falso, deja S3 en
Insuficiente.

\begin{tcolorbox}[nota]
\textbf{La declaración debe decir, con casos verificables:} en qué se usó IA
(análisis de código con agentes, generación de fichas, redacción, búsqueda de
literatura), cómo se verificó lo que devolvió (validaciones independientes con
múltiples agentes, \emph{y ejecución empírica real de los scripts}), y al menos
un caso concreto de algo \emph{descartado porque estaba mal} — el equipo ya
tiene tres documentados en sus propias validaciones: (1) un agente afirmó un
bug falso (el ``SÍ'' de \texttt{diversidad.py}) que se corrigió tras leer el
código; (2) un artefacto de codificación ``SÃ'' por ``Sí'' detectado y
corregido; (3) un orden de capas reportado como alfabético cuando era canónico
(artefacto del verificador). Además: riesgos éticos del proyecto (datos
personales y categorías sensibles: etnia, discapacidad, víctimas $\rightarrow$
Ley 1581 de 2012; sesgo de auto-declaración; riesgo de reidentificación) y el
\textbf{reparto real con nombres y horas}.
\end{tcolorbox}

\subsection{Requisitos de formato y citación}

\textbf{Documento de 6--10 págs.} (sin portada ni referencias), PDF, APA 7, con
la estructura: problema 1½--2 págs. (D1) · estado del arte 2--3 págs. (D2) ·
objetivos ½ pág. (D3) · datos y método 1--2 págs. (D4) · alcance/riesgos/
cronograma 1 pág. (D5) · declaración IA y reparto ½ pág. (S3) · referencias
APA 7. Los materiales actuales (informe base + auditoría de 126 págs. + SotA de
20 págs.) deben \textbf{sintetizarse}, no recortarse mecánicamente. La
conversión a APA 7 es mecánica pero indispensable: la bibliografía actual es
autor-año; ~18--25 referencias citadas bastan, todas rastreables.
